% Paper 3, §6 — Tree-search families: BFS, DFS, MCTS (anchor section)
% Anchors: kb/notes/reasoning/inference-time-search.md §1-3
%          kb/excerpts/rstar-math2025.md (#abstract, #headline, #sec-3, #sec-3-mcts)
%          kb/excerpts/snell2024.md (#sec-5-2-search)
%          (zhang2024-rest-mcts, wei2025-tree-search-survey, hu2025-mcts-rag
%           have no excerpt files yet; load-bearing claims tagged \deepencite.)

\section{Tree-search families: BFS, DFS, MCTS}
\label{sec:tree-search}

\Cref{sec:inference-compute-scaling} factored test-time compute into a
proposer leg and a verifier leg; \cref{sec:process-reward-models} built
out the verifier leg as a per-step scoring function $R^P_\phi$ over
partial prefixes $s_{1:t}$. The third decomposition --- after sampling
(\cref{sec:self-consistency}) and sequential revision
(\cref{sec:snell-tensor-shapes})~--- is \emph{tree search}: explicitly
construct and explore a tree of partial reasoning prefixes, using the
PRM as a value function to direct expansion. The taxonomy of the
2023--2025 literature, surveyed in
\citet{wei2025-tree-search-survey},\deepencite{wei2025-tree-search-survey
\S2--3 (taxonomy enumeration; excerpt file pending)} catalogues three
families that bracket the design space: beam over reasoning steps, BFS
or level-order expansion, and Monte-Carlo Tree Search. This section
formalises the search-tree datastructure, transcribes the canonical
reasoning-MCTS recurrence from
\citet{rstar-math2025},\footnote{KB excerpt:
\texttt{kb/excerpts/rstar-math2025.md\#sec-3-mcts}.} walks the four
phases of an MCTS iteration, situates ReST-MCTS\textsuperscript{*}
\citep{zhang2024-rest-mcts} as the AlphaZero-shaped reasoning loop,
notes MCTS-RAG \citep{hu2025-mcts-rag} as the retrieval-extended
cousin, and closes on the live contradiction between MCTS-headline
framings and \citet[\S 5--6]{snell2024}'s evidence that BFS-with-PRM is
competitive at fixed inference budget on math.

\subsection{The search tree as a datastructure over partial traces}
\label{sec:tree-datastructure}

Fix a problem $x$ and let $z_{1:t} = (s_1, \ldots, s_t)$ denote a
partial reasoning trace decomposed into $t$ steps, with each step
$s_i$ a contiguous span of trace tokens (a sentence, an equation
line, a bullet)%
\footnote{KB note:
\texttt{kb/notes/reasoning/inference-time-search.md\#1-formal-definition}.}.
The \textbf{search tree} has nodes indexed by partial traces $z_{1:t}$
and edges labelled by the next step $s_{t+1}$, with the root at
$z_{1:0} = \varnothing$. Three signals attach to the tree:
\begin{itemize}
\item the \emph{policy prior}
$\pi_\theta(s_{t+1} \mid x, z_{1:t})$, which assigns mass to candidate
next steps under the underlying LLM;
\item the \emph{value function}
$V_\phi(x, z_{1:t}) \to \mathbb{R}$, which scores partial states. The
canonical instantiation is $V_\phi = R^P_\phi$, the per-step PRM of
\cref{eq:prm-signature};
\item the \emph{terminal verifier}
$\mathcal{V}(x, y) \to \{0, 1\}$, which scores complete answers. For
math and code this is the same rule-based check that defines the RLVR
reward of \cref{eq:verifier-reward}.
\end{itemize}
The three search families differ entirely in how they combine
$\pi_\theta$ and $V_\phi$ to direct expansion; the datastructure is
common.

\paragraph{Symbol glossary.} For the rest of this section:
\begin{center}
\small
\begin{tabular}{ll}
\toprule
Symbol & Meaning \\
\midrule
$B$                 & beam width / branching factor at each expansion \\
$D$                 & maximum search depth (max number of steps) \\
$L_\text{step}$     & mean per-step token count \\
$L_\text{ans}$      & mean full-answer trace length, $L_\text{ans} \le D \cdot L_\text{step}$ \\
$N_\text{rollout}$  & number of MCTS simulation iterations \\
$N_s$               & MCTS visit count at node $s$ \\
$Q(s)$              & MCTS running mean rollout outcome through $s$ \\
$c_\text{puct}$     & UCT exploration constant \\
$V_\phi$            & learned value function on partial prefixes (typically a PRM) \\
$\mathcal{V}$       & terminal verifier (rule-based for math/code) \\
\bottomrule
\end{tabular}
\end{center}

\subsection{Beam search and BFS-V over reasoning steps}
\label{sec:tree-beam}

The simplest search family generalises decoder beam search from
token-level to step-level granularity. Maintain a frontier of $B$
partial traces; at each depth, expand each by sampling $k$ next-step
continuations from $\pi_\theta$; rescore the resulting $B \cdot k$
partials with $V_\phi$; retain the top $B$ by value. Iterating to
depth $D$ produces $B$ candidate full traces; the highest-$V_\phi$ or
highest-$\mathcal{V}$ trace is returned.

The cost is
\begin{equation}
  C_\text{beam} \;=\; O\!\left( B \cdot k \cdot D \cdot L_\text{step} \right)
  \quad\text{tokens.}
  \label{eq:cost-beam}
\end{equation}
\citet[\S 5.2]{snell2024} formalise the variant they call
\textbf{BFS-V}: a level-order beam with $N$ initial samples and beam
width $M$. The procedure (Snell's quoted text): ``sample $N$ initial
predictions for the first step in the solution; score the generated
steps according to the PRM's predicted step-wise reward-to-go estimate;
filter for only the top $\frac{N}{M}$ highest scoring steps; now from
each candidate, sample $M$ proposals from the next step, resulting in
a total of $N/M \times M$ candidate prefixes again. Then repeat steps
2--4''
\citep[\S 5.2]{snell2024}.\footnote{KB excerpt:
\texttt{kb/excerpts/snell2024.md\#sec-5-2-search}.} BFS-V differs from
classical beam only in the per-step scoring object: the PRM's
reward-to-go on $s_{1:t}$ replaces the language-model log-likelihood,
and the level-order structure means every depth-$t$ node is fully
evaluated before any depth-$(t+1)$ node is expanded.

\paragraph{Why beam-class methods are competitive.} Two structural
properties favour beam-shaped search on reasoning workloads. First,
the cost in \cref{eq:cost-beam} is linear in $D$ and in $B \cdot k$,
both directly tunable knobs that compose with the difficulty-bin
allocation of \cref{sec:snell-difficulty}: spend wider beams on harder
prompts, narrower on easier ones. Second, the per-step PRM signal is
\emph{exactly} what beam needs --- a scalar score on a partial prefix
--- and \cref{sec:prm-prm800k}'s min-aggregation result
($R^P_\text{min}(z) = \min_t R^P_\phi(x, s_{1:t})$) carries over node
by node: a beam node's score is the worst PRM score along its
root-to-node path, and the weakest-link inequality of
\cref{sec:prm-prm800k} prunes weak branches early. The DFS and
tree-of-thoughts variants
\deepencite{yao2023-tree-of-thoughts \S3 (canonical citation pending)}
add backtracking on a per-node ``promising / not promising''
heuristic, but the underlying datastructure is the same; they trade
the level-order regularity of BFS for depth-first traversal under a
node-evaluation gate.

\subsection{The MCTS recurrence (rStar-Math)}
\label{sec:tree-mcts}

Monte-Carlo Tree Search adds a fourth ingredient that beam lacks:
\emph{rollout-derived value estimates} are accumulated at each node and
used to bias future expansion via an exploration--exploitation rule.
\citet{rstar-math2025} take MCTS as the central inference-time
mechanism --- ``a math policy SLM performs test-time search guided by
an SLM-based process reward model''
\citep[abstract]{rstar-math2025}\footnote{KB excerpt:
\texttt{kb/excerpts/rstar-math2025.md\#abstract}.} --- and report 7B
policies reaching $90.0\%$ on MATH and $53.3\%$ on AIME 2024 after
four rounds of self-evolution
\citep[headline numbers]{rstar-math2025}.\footnote{KB excerpt:
\texttt{kb/excerpts/rstar-math2025.md\#headline}.}

\paragraph{The UCT selection rule.} At each iteration, descend the tree
from the root by repeatedly choosing the child $s$ that maximises%
\footnote{KB note:
\texttt{kb/notes/reasoning/inference-time-search.md\#2-2-mcts-over-reasoning-steps-rstar-math}.}
\begin{equation}
  U(s) \;=\; Q(s) \;+\; c_\text{puct}\;
              \pi_\theta\!\bigl(s \,\big|\, \text{parent}(s)\bigr)\;
              \frac{\sqrt{N_\text{parent}}}{1 + N_s}.
  \label{eq:uct}
\end{equation}
The first term $Q(s)$ is the running mean of rollout outcomes through
$s$ (the exploitation term: nodes that have produced correct rollouts
are preferred). The second term is the PUCT exploration bonus
\citep{rstar-math2025}\deepencite{rstar-math2025 \S3 (exact constants
pending PDF acquisition)}: the policy prior
$\pi_\theta(s \mid \text{parent})$ shapes which children are
\emph{a~priori} preferred, the
$\sqrt{N_\text{parent}}/(1 + N_s)$ ratio inflates as the parent has
been visited many times relative to $s$, and $c_\text{puct}$ dials the
overall exploration aggressiveness. The constant $c_\text{puct}$ is
the only free hyperparameter of the rule itself; rStar-Math sets it
empirically.

\paragraph{The four phases of an MCTS iteration.}
\begin{enumerate}
\item \textbf{Selection.} From the root, repeatedly apply
\cref{eq:uct} to descend until a leaf $\ell$ is reached. ``Leaf'' here
means a node that has not yet been expanded under MCTS bookkeeping ---
not necessarily a terminal in the trace sense.
\item \textbf{Expansion.} Sample $k$ candidate next-steps from
$\pi_\theta(\cdot \mid x, z_{\ell})$ and add them as children of
$\ell$, each with $Q = 0$ and $N = 0$.
\item \textbf{Simulation.} From one of the freshly-expanded children
$c$, roll out a complete trace under $\pi_\theta$ to a terminal state;
score the terminal answer with $\mathcal{V}$, and optionally weight
intermediate states by $V_\phi = R^P_\phi$ to obtain a real-valued
return $G_c \in [0, 1]$.
\item \textbf{Backup.} Propagate $G_c$ up the path from $c$ to root,
incrementing $N_s \mathrel{{+}{=}} 1$ and updating
$Q(s) \mathrel{{\leftarrow}} Q(s) + (G_c - Q(s))/N_s$ at each visited
node.
\end{enumerate}
After $N_\text{rollout}$ iterations, the answer is selected as either
the most-visited path from root, the rollout that achieved highest
return, or a PRM-weighted vote across rollouts; rStar-Math uses the
last\deepencite{rstar-math2025 \S3.x (answer-selection rule pending PDF)}.

\paragraph{Cost analysis.} Each MCTS iteration costs one rollout (a
full trace under $\pi_\theta$) plus $O(D)$ tree-traversal operations
plus $k$ expansion samples. The rollout dominates, so total cost is
\begin{equation}
  C_\text{MCTS} \;=\; O\!\left( N_\text{rollout} \cdot D \cdot L_\text{step} \right)
  \quad\text{tokens,}
  \label{eq:cost-mcts}
\end{equation}
matching the bound assumed by \cref{sec:snell-tensor-shapes}'s strategy
table. Comparing \cref{eq:cost-beam} and \cref{eq:cost-mcts}: at
matched depth $D$ and per-step length $L_\text{step}$, $N_\text{rollout}$
plays the role of $B \cdot k$. The two algorithms differ in
\emph{how} the budget is spent: beam fans out uniformly at each depth,
while MCTS allocates iterations adaptively to nodes whose $Q$-values
suggest the path is promising. On a benign loss landscape (good policy,
calibrated PRM), the adaptive allocation buys MCTS a constant-factor
edge; on rough landscapes (under-trained PRM, narrow correct path),
the same adaptivity can collapse onto a wrong branch.

\begin{intuition}
MCTS is a \emph{compute-allocator} that lets early rollout outcomes
re-shape where later rollouts go. Beam search, by contrast, allocates
the same width to every depth regardless of how the frontier evolves.
On problems where one branch has a sharp success advantage that
manifests early in rollouts, MCTS amplifies the signal --- $Q(s)$
rises, $N_s$ rises, the UCT bonus on alternate children shrinks, and
$N_\text{rollout}$ increasingly concentrates on the winning subtree.
On problems where rollout outcomes are noisy or where the success
signal is buried deep, MCTS spends iterations re-validating
already-committed paths instead of broadening, and beam can match or
beat it. Returning to math: \cref{eq:uct}'s exploitation term $Q(s)$
is a Monte-Carlo estimate of the expected return from $s$, and the
ratio $\sqrt{N_\text{parent}}/(1 + N_s)$ is the standard PUCT
exploration bonus from AlphaGo / AlphaZero, here applied in the
text-step action space rather than the board-position action space.
\end{intuition}

\subsection{Tensor-shape and operational notes}
\label{sec:tree-shapes}

The MCTS state is a node $s$ in the trace tree, indexed by its
root-to-node prefix $z_{1:t}$. The PRM is the value function on the
prefix: $V_\phi(s) = R^P_\phi(x, s_{1:t}) \in [0, 1]$, with the
aggregation rules of
\cref{eq:prm-min,eq:prm-prod,eq:prm-last} carrying over node by node.
Per-iteration cost decomposes as
\begin{center}
\small
\begin{tabular}{ll}
\toprule
Operation                          & Cost \\
\midrule
Selection (root $\to$ leaf $\ell$) & $O(D)$ tree-traversal ops, no LLM calls \\
Expansion ($k$ candidate children) & $k \cdot L_\text{step}$ tokens \\
Simulation (full rollout)          & $D \cdot L_\text{step}$ tokens, plus terminal $\mathcal{V}$ \\
Backup (root $\leftarrow$ leaf)    & $O(D)$ tree-traversal ops \\
\bottomrule
\end{tabular}
\end{center}
The \emph{rollout} dominates at typical $k \in [4, 16]$ and
$D \in [10, 50]$, which is why \cref{eq:cost-mcts} drops the expansion
and traversal terms. Operationally, the PRM call cost adds a constant
per simulation step --- $c_\text{verify}$ in
\cref{sec:snell-tensor-shapes}'s table --- but is normally amortised by
batched verifier inference across multiple rollouts in flight.

\subsection{ReST-MCTS\textsuperscript{*}: AlphaZero-shaped self-training}
\label{sec:tree-rest-mcts}

\citet{zhang2024-rest-mcts} elevate MCTS from an inference-time
mechanism to a \emph{training-data generator}.\deepencite{zhang2024-rest-mcts
\S3 (excerpt skeleton-only as of 2026-05; PDF acquisition pending)}
The protocol, drawn from the inference-time-search KB note's
treatment%
\footnote{KB note:
\texttt{kb/notes/reasoning/inference-time-search.md\#2-3-rest-mcts}.}:
\begin{enumerate}
\item Run MCTS with the current policy $\pi_\theta^{(r)}$ and PRM
$V_\phi^{(r)}$ to solve a batch of training problems.
\item Collect MCTS-discovered correct trajectories (those whose
terminal $\mathcal{V}(x, y) = 1$).
\item SFT the policy on these harvested trajectories:
$\pi_\theta^{(r+1)} \leftarrow \mathrm{SFT}\bigl(\pi_\theta^{(r)};
\{(x, z, y) : \mathcal{V}(x, y) = 1\}\bigr)$.
\item Update the PRM $V_\phi^{(r+1)}$ using the new policy's
trajectory distribution and the trajectory-correctness labels.
\item Repeat from step 1.
\end{enumerate}
The structural identity with AlphaZero
\deepencite{silver2017-alphazero \S algorithm (canonical citation
pending; KB note:
\texttt{kb/notes/reasoning/inference-time-search.md\#7-intuitions-and-analogies})}
is exact:
\begin{itemize}
\item self-play games $\leftrightarrow$ MCTS rollouts on math problems;
\item board positions $\leftrightarrow$ partial reasoning prefixes
$z_{1:t}$;
\item value head $\leftrightarrow$ PRM $V_\phi$;
\item policy network $\leftrightarrow$ next-step policy $\pi_\theta$;
\item self-play SFT $\leftrightarrow$ SFT on MCTS-discovered correct
trajectories.
\end{itemize}
rStar-Math \citep{rstar-math2025} runs a structurally similar loop
across four self-evolution rounds, with one important variation: the
PRM is replaced by a \emph{process preference model} (PPM) trained
from preferences over MCTS-discovered step-pairs rather than from
absolute step-level scores
\citep[\S 3]{rstar-math2025}.\footnote{KB excerpt:
\texttt{kb/excerpts/rstar-math2025.md\#sec-3}.} The choice trades
absolute calibration for a more robust pairwise signal that does not
require pinning the verifier's score scale; the structural role
inside the MCTS loop is unchanged.

\begin{analogy}
ReST-MCTS\textsuperscript{*} is AlphaZero \emph{with text steps}.
Selection, expansion, simulation, and backup are the same four phases.
The action space at each tree node is the LLM's next-step distribution
$\pi_\theta(s_{t+1} \mid x, z_{1:t})$ rather than the board's
legal-move distribution. The value head is the PRM
$R^P_\phi(x, s_{1:t})$ rather than the position-evaluation network.
The self-play loop is replaced by a
generate-MCTS-trajectories-then-SFT-on-correct-ones loop. Returning to
math: \cref{eq:uct}'s PUCT formula is the same equation that drives
AlphaZero's selection, with the policy prior
$\pi_\theta(s \mid \text{parent})$ playing the role of the
move-prior network's output and $Q(s)$ playing the role of the
running-mean win rate. The analogy is load-bearing because it
explains why the iterative SFT-on-search-discovered-correct loop
converges: each round's MCTS uses a stronger policy and a stronger
value head than the last, and the harvested trajectories are
strictly better-than-policy on average, supplying a positive-gradient
SFT signal that ratchets the policy upward without ever introducing
a separate value-network bootstrap.
\end{analogy}

\subsection{MCTS-RAG: tree search over retrieval and reasoning}
\label{sec:tree-mcts-rag}

\citet{hu2025-mcts-rag} extend the MCTS state space to include
\emph{retrieval} actions alongside reasoning
steps.\deepencite{hu2025-mcts-rag \S3 (excerpt skeleton-only as of
2026-05; PDF acquisition pending)} Inside an MCTS iteration, the
expansion step samples either a next reasoning step from
$\pi_\theta(\cdot \mid x, z_{1:t})$ or a retrieval call into a
document index, with the choice itself an action under a hybrid
policy. The value function $V_\phi$ is generalised to score states
that mix reasoning and retrieved evidence; the terminal verifier
$\mathcal{V}$ remains domain-specific (a fact-match check on a
knowledge-grounded QA benchmark, an arithmetic check on a
calculation-grounded one).

The structural advance: tree search lets the model decide
\emph{when} to retrieve and \emph{what} to retrieve as part of the
reasoning trajectory, rather than fixing retrieval as a one-shot
preprocessing step. The cost analysis of \cref{eq:cost-mcts}
generalises: the per-step token cost $L_\text{step}$ is now a mixture
over reasoning-step length and retrieval-call length (typically
shorter, but with an additional index-lookup latency that the cost
expression does not capture). MCTS-RAG closes the gap between two
otherwise-disjoint inference-time-compute literatures --- search-based
TTC for reasoning and RAG for knowledge access --- under one
tree-search framework.

\subsection{Search-vs-trained-CoT: the live contradiction}
\label{sec:tree-search-vs-cot}

The 2024 hypothesis articulated by the early reasoning-MCTS literature
was that explicit search at inference time was \emph{necessary} for
strong reasoning performance: a small policy plus heavy MCTS was
expected to dominate a large policy without search. The 2025--2026
evidence is more nuanced. \citet[\S 4]{rstar-math2025} demonstrate
the small-policy-with-heavy-MCTS extreme: a 7B policy plus four rounds
of MCTS self-evolution reaches 90.0\% on MATH and 53.3\% on AIME 2024,
matching or exceeding o1-preview at much smaller policy
scale.\footnote{KB excerpt:
\texttt{kb/excerpts/rstar-math2025.md\#headline}.} \citet{deepseek-r1}
demonstrate the opposite extreme: a large policy with RLVR-trained
long CoT and \emph{no} explicit inference-time search reaches
77.9\% pass@1 on AIME 2024 and 86.7\% with self-consistency, on the
strength of trace-length growth from $\sim$200 to $\sim$10\,000 tokens
during training (\cref{sec:r1-protocol}). The two extremes occupy
roughly the same accuracy region at very different policy scales and
inference-time-compute profiles.

The clean reading%
\footnote{KB note:
\texttt{kb/notes/reasoning/inference-time-search.md\#5-the-search-vs-trained-cot-debate}.}
is that explicit search is most valuable where the policy is weak
(small model, narrow domain, hard tail). At strong policy plus
standard distributions, internalised search via long CoT is
competitive at matched compute. This matches the difficulty-bin
operationalisation of \cref{sec:snell-difficulty}: search wins on
the hard bin, sampling/refinement wins on the easy bin, and both
collapse to the same answer on prompts the policy nearly solves.

\begin{contradiction}
\textbf{BFS-with-PRM vs.\ MCTS at fixed inference budget.}
\citet[\S 5--6]{snell2024} report BFS-V (the level-order
beam-with-PRM of \cref{sec:tree-beam}) competitive with or outright
better than alternative search strategies at matched compute on the
MATH benchmark with PaLM~2-S\* as the base model
\citep[\S 5.2]{snell2024}.\footnote{KB excerpt:
\texttt{kb/excerpts/snell2024.md\#sec-5-2-search}.} The implication
is that the extra MCTS machinery (UCT bookkeeping, rollout-derived
$Q(s)$ updates, PUCT exploration constant $c_\text{puct}$) does not
necessarily pay for itself at fixed test-time FLOPs. \citet{rstar-math2025}'s
headline framing instead foregrounds MCTS as the central mechanism by
which 7B SLMs match o1-preview, and the AlphaZero structural parallel
of \cref{sec:tree-rest-mcts} runs entirely through the MCTS recurrence
of \cref{eq:uct}. Reconciling the two: Snell's evaluation uses a
non-reasoning-trained base model and a single round of search; rStar-
Math's MCTS is co-trained with the policy and PRM across four rounds
of self-evolution \citep[abstract]{rstar-math2025}, so the
``MCTS'' that wins is a different object from the
``MCTS'' that ties. Whether the MCTS-vs-BFS gap reverses on
reasoning-trained base models, on harder benchmarks, or under
multi-round co-evolution is not closed in 2026; the
\citet{wei2025-tree-search-survey} taxonomy
catalogues the disagreement without resolving
it.\deepencite{wei2025-tree-search-survey \S5 (resolution discussion;
excerpt file pending)} \Cref{sec:search-vs-rl} returns to the
search-vs-RL face of the same tradeoff.
\end{contradiction}

\subsection{Forward link}
\label{sec:tree-forward}

Tree search is the verifier-leg's most expressive realisation. Its
cost-per-trajectory dominates self-consistency
(\cref{sec:self-consistency}) and best-of-$N$
(\cref{sec:prm-prm800k}), and its operational complexity dominates
beam search; the question is whether it earns its budget. The answer,
across \cref{sec:tree-mcts}--\cref{sec:tree-search-vs-cot}, depends
on the policy's strength, the value function's calibration, and
whether the search is co-evolved with the policy
(rStar-Math, ReST-MCTS\textsuperscript{*}) or run once over a fixed
checkpoint (Snell-style evaluation).

\Cref{sec:rlvr-grpo-dapo} introduces the third leg of the reasoning
triangle --- RL on verifiable rewards --- which is the
\emph{training-time} answer to the question this section treated at
inference: how to spend compute to amplify reasoning. The
\cref{eq:cost-mcts} cost expression and the \cref{eq:uct} UCT
recurrence both reappear there in a different role: ReST-MCTS\textsuperscript{*}'s
search-then-SFT loop and RLVR's group-of-rollouts both \emph{generate}
training trajectories from a current policy, with the verifier
$\mathcal{V}$ deciding which trajectories survive. \Cref{sec:search-vs-rl}
then puts curves on the tradeoff, asking at fixed total compute when
marginal FLOPs are best spent on more inference search
(\cref{sec:tree-mcts}), more RLVR training (\cref{sec:rlvr-grpo-dapo}),
or more pretraining (\cref{sec:snell-flops-matched}). The MCTS-vs-BFS
contradiction of \cref{sec:tree-search-vs-cot} and the absent
closed-form $L(\theta, N)$ flagged in \cref{sec:snell-open} are two
faces of the same open frontier: the search-leg's quantitative
contribution to inference-time-compute returns is precisely what
remains unresolved in 2026.
